\makeatletter
\def\input@path{{./}{../}{../../}{../../../}{../../../../}{preamble/}{../preamble/}{../../preamble/}{../../../preamble/}}
\makeatother
\input{preamble_exams.tex}
\graphicspath{{./}{../}{../../}{../../../}{../../../../}{preamble/}{../preamble/}{../../preamble/}{../../../preamble/}}
\setYColumnWidth{0.30}
% Keep the inherited 11pt exam typography; allow for the heading box padding.
\let\ExamSubsectionBox\SubsectionBox
\renewcommand{\SubsectionBox}[1]{\begingroup
\setlength{\linewidth}{\dimexpr\linewidth-2\fboxsep\relax}
\ExamSubsectionBox{#1}\endgroup}

\begin{document}
\ExamHeader{10}{3}{DECISION-MAKING UNDER UNCERTAINTY}{1st}
\vspace{1mm}
\LearningObjective{Compare decision criteria and their limitations, apply staged filtering, and plan a justified investigation of a real-life economic or financial decision under uncertainty.}{
C8: Compares the Maximax, Maximin, Minimax Regret, and EMV/Bayes decision criteria by contrasting their information requirements, risk orientations, recommended alternatives, and limitations.\par
C9: Constructs and interprets decision trees for staged filtering rules, applying decision criteria in the specified order to identify the final alternative.\par
C10: Evaluates economic and financial decision-making scenarios by selecting and justifying the most appropriate decision criterion according to the information available, risk attitude, and decision context.}

\StudentInfo

\begin{multicols}{2}
\raggedcolumns
\SubsectionBox{Criteria assessment}
Answer individually on separate paper. Questions 1--6 are worth 2 points each; Question 7 has three components worth 2 points each. A \textbf{minimum of 4 points} is required to pass each criterion. Partial credit is available for correct comparisons, filtering stages, and justified reasoning even if later conclusions contain errors.

\begin{center}
\begin{tabular}{|c|c|c|c|}\hline
\textbf{Question}&\textbf{C8}&\textbf{C9}&\textbf{C10}\\\hline
1&2&--&--\\\hline
2&2&--&--\\\hline
3&2&--&--\\\hline
4&--&2&--\\\hline
5&--&2&--\\\hline
6&--&2&--\\\hline
7(a), (b), (c)&--&--&$2+2+2$\\\hline
\textbf{Total}&\textbf{6}&\textbf{6}&\textbf{6}\\\hline
\end{tabular}
\end{center}



\SubsectionBox{1. Repair Workshop Plans \score{C8-\sml{2}}}
A bicycle repair workshop chooses its capacity for the next quarter before demand is known. Profits are in thousand USD; forecast probabilities are shown below. \vspace{-3mm}
\begin{center}\small\setlength{\tabcolsep}{4pt}
	\begin{tabular}{lrrr}\toprule
		Alternative & High & Normal & Low\\
		$P(S_j)$ & 0.2 & 0.6 & 0.2\\ \midrule
		Large ($A_1$) & 120 & 25 & -20\\
		Flexible ($A_2$) & 85 & 65 & 10\\
		Basic ($A_3$) & 50 & 45 & 35\\
		Partner ($A_4$) & 70 & 75 & 20\\
		\bottomrule\end{tabular}\end{center}

The following table summarizes the decision-making statistics for each alternative.\vspace{-3mm}

\begin{center}\small\setlength{\tabcolsep}{4pt}
	\begin{tabular}{lrrrr}\toprule
		& Max & Min & Regret & EMV\\ \midrule
		$A_1$ & 120 & -20 & 55 & 35\\
		$A_2$ & 85 & 10 & 35 & 58\\
		$A_3$ & 50 & 35 & 70 & 44\\
		$A_4$ & 75 & 20 & 50 & 63\\
		\bottomrule\end{tabular}\end{center}
	\vspace{-3mm}

\begin{enumerate}[label=\textbf{(\alph*)},leftmargin=*]
	\item Report each criterion's recommended alternative(s), deciding statistic, and value. \score{1}
	
	\item A subsidy $s$ (integer, $0\leq s\leq30$, thousand USD) increases the low-demand payoffs of both $A_1$ and $A_2$ by $s$. The rest remain unchanged. An alternative is eligible only if its profit is nonnegative in every state.	
	Determine the ranges of $s$ for which each alternative is eligible. Find the smallest $s$ making $A_2$ the unique Maximin choice. At that $s$, report all four recommendations and deciding values among eligible alternatives. \score{1}
\end{enumerate}
\end{multicols} 	

\newpage
\begin{multicols}{2}
\raggedcolumns
\SubsectionBox{2. School Store Orders \score{C8-\sml{2}}}
A school store chooses one stationery order for the coming term. Net profits are in hundred USD.
\begin{center}\small\setlength{\tabcolsep}{4pt}
\begin{tabular}{lrrr}\toprule
Alternative & High & Normal & Low\\
$P(S_j)$ & 0.25 & 0.5 & 0.25\\ \midrule
Order $B_1$ & 72 & 54 & 36\\
Order $B_2$ & 60 & 48 & 24\\
Order $B_3$ & 72 & 42 & 36\\
\bottomrule\end{tabular}\end{center}
\begin{center}\small\setlength{\tabcolsep}{4pt}
\begin{tabular}{lrrrr}\toprule
& Max & Min & Regret & EMV\\ \midrule
$B_1$ & 72 & 36 & 0 & 54\\
$B_2$ & 60 & 24 & 12 & 45\\
$B_3$ & 72 & 36 & 12 & 48\\
\bottomrule\end{tabular}\end{center}
\begin{enumerate}[label=\textbf{(\alph*)},leftmargin=*]
\item Report all four recommendations and deciding values, including every tie. \score{1}
\item A supplier charges an extra flat fee $f$ for $B_1$, deducted once from its profit in every state. Here $f$ is an integer, $0\leq f\leq10$, in hundred USD. All other payoffs and probabilities stay fixed; all orders remain eligible. Find the smallest $f$ making EMV/Bayes strictly prefer $B_3$ over $B_1$. Justify the threshold, identify the boundary tie, and interpret the fee in USD. At that $f$, report all four updated recommendations and deciding values, including every tie. \score{1}
\end{enumerate}

\columnbreak
\SubsectionBox{3. Community Market Stall \score{C8-\sml{2}}}
An organizer chooses a stall for one weekend. Profits are in hundred USD. The forecast assigns high attendance probability 0.40 and low attendance probability 0.60.
\begin{center}\small\setlength{\tabcolsep}{4pt}
\begin{tabular}{lrr}\toprule
Alternative & High & Low\\
$P(S_j)$ & 0.4 & 0.6\\ \midrule
Large ($C_1$) & 80 & 10\\
Compact ($C_2$) & 48 & 32\\
\bottomrule\end{tabular}\end{center}
\begin{center}\small\setlength{\tabcolsep}{4pt}
\begin{tabular}{lrrrr}\toprule
& Max & Min & Regret & EMV\\ \midrule
$C_1$ & 80 & 10 & 22 & 38\\
$C_2$ & 48 & 32 & 32 & 38.4\\
\bottomrule\end{tabular}\end{center}
\begin{enumerate}[label=\textbf{(\alph*)},leftmargin=*]
\item For part (a) only, the forecast is withdrawn as unreliable. Compare the information needed by all four criteria: which summaries still support a choice, which no longer do, and why would equal probabilities need justification? \score{1}
\item For this part, restore the supplied probabilities as reliable. A sponsor guarantees $C_2$ a high-attendance profit of $g$, replacing 48, where $g$ is an integer, $48\leq g\leq80$, in hundred USD. All other payoffs stay fixed; both stalls remain eligible. Find the smallest $g$ making $C_2$ the unique Minimax Regret choice. Justify the threshold and boundary tie; interpret the increase from 48 in USD. At that $g$, report all four updated recommendations and deciding values, including every tie. \score{1}
\end{enumerate}

\end{multicols}
\newpage
\begin{multicols}{2}
\raggedcolumns
\SubsectionBox{Staged filtering: Questions 4--6}
For \textbf{each} problem, construct a compact filtering tree with branches for retained and eliminated alternatives. Label every stage with its rule and relevant numerical values. Identify the final choice, including any tie, and write one sentence interpreting the result in context. \score{1}

Also construct a supporting table with the following columns, recording every stage and its relevant values. \score{1}
\begin{center}\footnotesize\setlength{\tabcolsep}{2pt}
\begin{tabular}{@{}p{0.11\linewidth}p{0.23\linewidth}p{0.28\linewidth}p{0.28\linewidth}@{}}
\toprule
\# Step & Description & Left Alternatives & Eliminated Alternatives\\
\bottomrule
\end{tabular}\end{center}
\begin{itemize}[leftmargin=*]
\item Step 0 includes all six alternatives.
\item Each stage compares only survivors from the previous stage. Eliminated alternatives cannot reenter.
\item Supplied summary statistics remain fixed throughout filtering. Do not recalculate regrets or other statistics after eliminations.
\item When keeping the best $k$ alternatives, retain all tied at the cutoff, even if more than $k$ survive.
\item Report any final tie explicitly; do not invent a tie-breaker.
\end{itemize}
Use highest EMV for EMV/Bayes, highest Max for Maximax, highest Min for Maximin, and lowest Regret for Minimax Regret. Regrets were calculated against all six original alternatives.

\columnbreak
\SubsectionBox{4. Book Fair Stock \score{C9-\sml{2}}}
A bookseller chooses a stock plan $D_1$--$D_6$ for a fair. Supplied summaries describe net profits in hundred USD under high, normal, and low sales, with probabilities 0.50, 0.30, and 0.20.
\begin{center}\small\setlength{\tabcolsep}{4pt}
\begin{tabular}{lrrrr}\toprule
& EMV & Max & Min & Regret\\ \midrule
$D_1$ & 65 & 100 & 0 & 40\\
$D_2$ & 65 & 90 & 10 & 30\\
$D_3$ & 65 & 80 & 20 & 20\\
$D_4$ & 59 & 70 & 30 & 30\\
$D_5$ & 53 & 60 & 40 & 40\\
$D_6$ & 43 & 50 & 30 & 50\\
\bottomrule\end{tabular}\end{center}
Apply these stages in order:
\begin{enumerate}[leftmargin=*]
\item EMV/Bayes: keep the 2 highest EMVs, retaining cutoff ties.
\item Maximin: select the highest Min among survivors.
\end{enumerate}
Complete the tree, table, and interpretation described above.
\SubsectionBox{5. Music Event Packages \score{C9-\sml{2}}}
An event organizer chooses one equipment package $E_1$--$E_6$ for a weekend music event. Summaries show net profits in hundred USD under high, normal, and low ticket sales, with probabilities 0.40, 0.40, and 0.20.
\begin{center}\small\setlength{\tabcolsep}{4pt}
\begin{tabular}{lrrrr}\toprule
& EMV & Max & Min & Regret\\ \midrule
$E_1$ & 58 & 110 & -10 & 55\\
$E_2$ & 63 & 95 & 15 & 30\\
$E_3$ & 69 & 95 & 25 & 20\\
$E_4$ & 67 & 80 & 35 & 30\\
$E_5$ & 61 & 70 & 45 & 40\\
$E_6$ & 52 & 60 & 40 & 50\\
\bottomrule\end{tabular}\end{center}
Apply these stages in order:
\begin{enumerate}[leftmargin=*]
\item Maximax: keep the 3 highest Max values.
\item Maximin: keep the 2 highest Min values among survivors.
\item Minimax Regret: select the lowest Regret among survivors.
\end{enumerate}
Construct the filtering tree and supporting table; include rules, values, eliminations, final choice, and a one-sentence interpretation. Use the shared rules for Questions 4--6.

\end{multicols}
\newpage
\begin{multicols}{2}
\raggedcolumns
\SubsectionBox{6. Refill Shop Expansion \score{C9-\sml{2}}}
A refill shop chooses an expansion plan $F_1$--$F_6$ for the next year. Summaries show net profits in thousand USD under strong, steady, and weak demand, with probabilities 0.30, 0.50, and 0.20.
\begin{center}\small\setlength{\tabcolsep}{4pt}
\begin{tabular}{lrrrr}\toprule
& EMV & Max & Min & Regret\\ \midrule
$F_1$ & 61 & 120 & 0 & 45\\
$F_2$ & 64.5 & 105 & 15 & 30\\
$F_3$ & 68.5 & 95 & 25 & 25\\
$F_4$ & 70 & 85 & 35 & 35\\
$F_5$ & 64 & 75 & 45 & 45\\
$F_6$ & 55 & 65 & 40 & 55\\
\bottomrule\end{tabular}\end{center}
Apply these stages in order:
\begin{enumerate}[leftmargin=*]
\item Minimax Regret: keep the 5 lowest Regret values.
\item EMV/Bayes: keep the 4 highest EMVs among survivors.
\item Maximax: keep the 2 highest Max values among survivors.
\item Maximin: select the highest Min among survivors.
\end{enumerate}
Construct the filtering tree and supporting table; include rules, values, eliminations, final choice, and a one-sentence interpretation. Use the shared rules for Questions 4--6.

\columnbreak
\SubsectionBox{7. Future Research Proposal \score{C10-\sml{6}}}
Write an \textbf{individual proposal} for the real-life economic or financial decision under uncertainty you intend to investigate for the C10 practice assignment. Organize your response into the three components below, using short, connected paragraphs.

This exam asks for your \textbf{research plan only}. You do not need completed research, collected data, a numerical payoff table, calculations, a final recommended alternative, completed references, or a PDF article. You may name intended sources; do not fabricate data or citations.

\textbf{(a) Proposed decision and context.} Describe a specific, realistic situation, the decision-maker, objective, and time period. Identify at least two feasible alternatives and at least two mutually exclusive states of nature covering the relevant possibilities. Explain the relevant risk attitude or constraints, such as budget or acceptable loss. \score{2}

\textbf{(b) Planned evidence and financial model.} Explain what financial information you will collect and where you expect to obtain it. Describe how revenues, costs, net gains, or savings will be used to construct a comparable payoff table in which larger values are better. State the intended currency, units, and time period. Explain how you will distinguish sourced facts from assumptions and estimates. \score{2}

\textbf{(c) Proposed criterion and justification.} Make a provisional choice among Maximax, Maximin, Minimax Regret, and EMV/Bayes. Explain why it fits the anticipated information, objective, risk attitude, and constraints. Briefly explain why each of the other three may be less suitable. \score{2}

If you propose EMV/Bayes, explain how you will obtain defensible probabilities for all states. Otherwise, explain why a criterion without probabilities fits. Identify one expected limitation or assumption that could affect the eventual recommendation.

Your criterion choice is \textbf{provisional} and may change after gathering evidence. The plan will be assessed for specificity, coherence, feasibility, and justification.
\end{multicols}
\end{document}
